\documentclass[11pt,a4paper]{article}
\usepackage[utf8]{inputenc}
\usepackage[T1]{fontenc}
\usepackage[english]{babel}
\usepackage{amsmath,amssymb}
\usepackage[round]{natbib}
\usepackage{booktabs}
\usepackage{graphicx}
\usepackage[margin=2.5cm]{geometry}
\usepackage{hyperref}
\usepackage{xcolor}
\hypersetup{colorlinks=true,linkcolor=blue!50!black,citecolor=blue!50!black,urlcolor=blue!50!black}
\usepackage{caption}
\captionsetup{font=small,labelfont=bf}

\title{\bfseries The Query Cannot See the Question\\
\large A Short Convolution's Reach Decides Which Part of a Query Conditions Retrieval,\\
and What Falls Outside Becomes a Confident Wrong Answer}

\author{Maximiliano Speranza\\
\small Independent Researcher, Buenos Aires, Argentina\\
\small \href{https://orcid.org/0009-0005-0413-8554}{ORCID 0009-0005-0413-8554}\\
\small \texttt{maximiliano.speranza@gmail.com}}

\date{September 2026}

\begin{document}
\maketitle

\begin{abstract}
\noindent
Linear-attention and state-space models form their query with a short causal depthwise convolution.
Kernel size $4$ is the de facto default, so the query at a given layer is a function of the current
token and three predecessors. We show that this reach is not a detail of local smoothing but a hard limit on
\emph{which part of a question is allowed to condition retrieval}, and that what falls outside does
not degrade gracefully: it disappears exactly, and the model answers confidently from the part it can
still see.

We measure the effect in a small language model with a co-trained persistent archive queried across
sequences. The sensitivity of retrieval to a query token, measured as the total-variation movement of
the read distribution when that token is replaced, is $\mathbf{0.000000}$ as soon as the token passes
the convolution's reach, and the cut-off moves with the kernel: with kernel $3$ (reach $2$) the step
falls between distances $2$ and $3$, and with kernel $5$ (reach $4$) between $4$ and $5$. Sixty cells
out of sixty, two architectures, three seeds each, two query components, five distances.

The behavioural consequence is a specific and common failure. In our task, questions have the form
\emph{what is the \textbf{<relation>} of \textbf{<entity>}?}, with the entity at distance $1$ from the
read position and the relation at distance $3$. With kernel $3$, the relation is outside the window in
$100\%$ of queries, so the model retrieves on the entity alone. When asked about a relation that was
never stated for an entity that \emph{was}, it answers from the wrong stored fact instead of
abstaining: correct abstention is $0.5850$ to $0.7349$ across three seeds. Widening the kernel to $5$,
which costs $1{,}280$ parameters out of $865{,}395$, lifts it to $\mathbf{0.9931}$ to
$\mathbf{1.0000}$ with no overlap between conditions, and overall accuracy to $0.988$--$0.993$ against
a trivial floor of $0.4065$.

We then check the mechanism outside our own model. In \texttt{mamba-130m-hf}, $129$M parameters,
changing one token five to eight positions before the read moves the layer output at every distance
and leaves the convolution output at \textbf{exactly zero}, $80$ cells out of $80$. The state sees the
whole sequence; the query that reads it does not. Incidentally, the oldest convolution tap in that
checkpoint is exactly zero in all $24$ layers, so its effective reach is $2$ rather than the nominal
$3$; we report the measurement and not a cause.

We give the diagnosis as a recipe that requires no training: change one token of the query and see
whether retrieval moves. If it does not, that token is invisible to the search, and any confidence the
model expresses about it is unearned.
\end{abstract}

\section{The problem}

A model that keeps a persistent memory has to do two things that are usually measured together and are
in fact separate. It has to \emph{find} the right entry, and it has to \emph{know} when the entry it is
looking for is not there. The second one is what makes the difference between a memory and a
confabulator, and it is the harder of the two.

We found that in our own system the second failure had a purely mechanical cause, upstream of anything
the model learns. The query used to search the archive is formed by a short causal convolution over the
hidden states. In a question of the form \emph{what is the \texttt{<art>} \texttt{<noun>} of
\texttt{<entity>}?}, the entity sits one token before the read position and the noun that names the
relation sits three tokens before it. A kernel of size $3$ reaches two tokens back. The relation is one
token outside the window, deterministically, in every single query.

The model was not ignoring the relation. It could not see it.

\section{Related work}

\textbf{Short convolutions are ubiquitous and their size is inherited, not measured.} Mamba-2
\citep{mamba2} uses a causal depthwise convolution of size $4$ by default. In the reference
implementation of Qwen3-Next, which uses Gated DeltaNet \citep{gateddeltanet} for its linear-attention
layers, the configuration exposes \texttt{linear\_conv\_kernel\_dim}, documented as ``kernel size of the
convolution used in linear attention layers'' and \textbf{defaulting to 4} \citep{qwen3next}. Ablations
in this line report that removing the convolution hurts, which establishes that local context matters;
none reports what is lost when the window is present but too narrow for the query at hand. A reach of
three tokens is therefore not an exotic configuration, it is the setting in deployed models.

\textbf{The closest theoretical result.} \citet{cat} study $N$-gram associative recall and prove, by
construction, that a causal filter of length $N$ suffices, matching the filter to the $n$-gram. That is
the same family of idea and we cite it as the precursor. Three things separate it from what we report.
Their query is a \emph{contiguous} $n$-gram, the last $N$ tokens with nothing irrelevant in between,
whereas what governs our case is the \emph{distance} of a component with filler between the two parts
that matter. Their Theorem 1 is an existence result and they run no experiment in which the filter is
too short. And they report the opposite in practice, stating that ``while K/Q convolution helps in
theoretical constructions for N-gram AR, in real experiments, they don't provide noticeable performance
benefits''. We measure a large effect in a regime they do not test, namely a persistent
\emph{cross-sequence} memory with a compositional query.

\textbf{Why the field could not see it.} The canonical benchmark for this capability, MQAR
\citep{zoology}, has a \emph{single-token} query: the key is looked up and the value returned. By
construction such a query cannot have one part inside the window and another part outside. Zoology
attributes $82\%$ of the recall gap to the convolution applying a fixed, input-independent filter, which
is a statement about \emph{adaptivity}. Ours is about \emph{reach}, and the two are independent: our
kernel-$5$ model has exactly the same input-independence and the failure disappears.

\textbf{From the attention side.} Multi-token attention \citep{mta} convolves over attention weights so
that attention can be conditioned on several tokens at once, motivated by the observation that a single
token is not enough to locate a match. It moves in the same direction and confirms ours by contrast: it
never asks which part of the query goes blind, and does not connect the question to abstention.

\section{Setup}

The substrate is a $3.5$\,MB language model, $865{,}395$ parameters, $d=128$, four blocks, trained from
scratch on a closed synthetic language of $242$ tokens with a co-trained persistent archive. Facts are
stated in one sequence, revised in another and queried in a third, \textbf{with the recurrent state
reset between sequences}, so retrieval cannot be solved by carrying activations forward.

Queries are read at the final token of the question. The read is injected in block $0$
\textbf{before the mixer}, which matters for the mechanism: at that point the hidden state is still the
token embedding, with no recurrent mixing, so the convolution's window is literally everything the
query can see. This also explains why the effect below can be exactly zero rather than merely small.

Four answer categories are scored separately. \texttt{current} and \texttt{previous} are questions with
an answer in the archive. \texttt{nose\_ent} asks about an entity never mentioned, the easy absence.
\texttt{nose\_rel} asks about a relation never stated \emph{for an entity that was}, which is the case
that looks like a real hallucination, because the retrieval finds a genuine neighbouring fact.

\section{Result 1: the window is a step function, and it is exactly zero outside}

We measure sensitivity by replacing one component of the question with a different one of the same
type and recording the total-variation distance between the read distributions. To vary distance
without touching the question, the read position is advanced $r$ places over the padding that already
follows the question mark, which is exactly equivalent to appending $r$ fillers.

\begin{center}
\begin{tabular}{lccccccc}
\toprule
model & kernel (reach) & $d{=}1$ & $d{=}2$ & $d{=}3$ & $d{=}4$ & $d{=}5$ & $d{=}6$ \\
\midrule
\texttt{v3} (3 seeds) & 3 (2) & $0.71$--$0.74$ & $0.15$--$0.25$ & $\mathbf{0.000000}$ & $\mathbf{0.000000}$ & $\mathbf{0.000000}$ & $\mathbf{0.000000}$ \\
\texttt{kq3} (3 seeds) & 5 (4) & $0.36$--$0.46$ & $0.10$--$0.15$ & $0.05$--$0.35$ & $0.07$--$0.33$ & $\mathbf{0.000000}$ & $\mathbf{0.000000}$ \\
\bottomrule
\end{tabular}
\end{center}

Sixty cells out of sixty match the prediction made before running, namely that sensitivity is positive
for $d \le$ reach and zero beyond it. One cell initially read $0.0106$ where a zero belonged; the cause
was that the ``previous value'' phrasing is two tokens longer, so advancing the read position clipped
against the tensor bound, and clipping \emph{shortens} the distance being measured. Those samples are
now discarded rather than clipped and the cell reads $0.000000$.

\section{Result 2: widening the window buys abstention}

\begin{center}
\begin{tabular}{lccccc}
\toprule
unit & kernel & \texttt{current} & \texttt{nose\_ent} & \textbf{\texttt{nose\_rel}} & false abstention \\
\midrule
\texttt{kq3\_s0} & 5 & $1.0000$ & $0.9538$ & $\mathbf{0.9931}$ & $0.0000$ \\
\texttt{kq3\_s1} & 5 & $0.9964$ & $0.9726$ & $\mathbf{1.0000}$ & $0.0030$ \\
\texttt{kq3\_s2} & 5 & $1.0000$ & $0.9356$ & $\mathbf{1.0000}$ & $0.0000$ \\
\midrule
\texttt{v3\_s0} & 3 & $1.0000$ & $0.9851$ & $0.6090$ & $0.0000$ \\
\texttt{v3\_s1} & 3 & $1.0000$ & $1.0000$ & $0.5850$ & $0.0000$ \\
\texttt{v3\_s2} & 3 & $0.9927$ & $0.9414$ & $0.7349$ & $0.0064$ \\
\bottomrule
\end{tabular}
\end{center}

There is no overlap between conditions on \texttt{nose\_rel}. Overall accuracy, counting a correct
abstention as correct, is $0.988$--$0.993$ for kernel $5$ against a trivial floor of $0.4065$, the score
of a model that abstains on everything.

\textbf{The cost is real and we state it.} \texttt{nose\_ent}, the easy case, drops slightly, from
$0.941$--$1.000$ to $0.936$--$0.973$. A reading that only reports a pooled abstention number hides this
trade.

\textbf{It is not capacity.} The kernel-$5$ model adds $1{,}280$ parameters, $0.15\%$, and the kernel-$3$
control already retrieves the right entry with probability $1.0000$. What changes is access, and Result
1 measures the access directly.

\section{Result 3: ablating the tap that reaches the relation}

Turning off one convolution tap at a time in the trained kernel-$5$ model, without retraining, damages
\texttt{current} by $0.483$, $0.582$ and $0.530$ for the tap that reaches the relation, against $0.372$,
$0.013$ and $0.195$ for the tap that reaches a function word and $0.218$, $0.093$ and $0.157$ for the
tap that reaches the article. The relation tap has the \emph{smallest} learned magnitude of the five,
so per unit of weight it costs about twice as much as either control.

A pre-registered criterion phrased on the abstention metric failed, and we report why: ablating the tap
removes the relation from \emph{every} query, so the model abstains more, and a metric that rewards
abstention cannot fall. The damage appears in accuracy and in false abstention instead.

\section{Result 4: the dissociation holds in a real model}

The results above are measured in a small model, so we checked the mechanism in one that is not ours:
\texttt{state-spaces/mamba-130m-hf}, $129{,}135{,}360$ parameters, on CPU, with no training. In Mamba
the short convolution is applied to $x$ \emph{before} $B$, $C$ and $\Delta$ are computed, and $C_t$ is
the analogue of the read query.

We change a single token at distance $d$ from a read position and measure, in layer $0$, the movement
of the convolution output and of the layer output. Ten read positions, two texts.

\begin{center}
\begin{tabular}{lcccccc}
\toprule
& $d{=}1$ & $d{=}2$ & $d{=}3$ & $d{=}5$ & $d{=}7$ & $d{=}8$ \\
\midrule
\texttt{conv1d} output & $0.8$--$3.7$ & $0.5$--$1.6$ & $\mathbf{0.0}$ & $\mathbf{0.0}$ & $\mathbf{0.0}$ & $\mathbf{0.0}$ \\
layer output & $9\!\times\!10^{-2}$ & $2\!\times\!10^{-2}$ & $2\!\times\!10^{-2}$ & $1\!\times\!10^{-2}$ & $9\!\times\!10^{-2}$ & $1\!\times\!10^{-2}$ \\
\bottomrule
\end{tabular}
\end{center}

The layer output moves at \emph{every} distance, $80$ cells out of $80$: the state does see the whole
sequence. The convolution output is exactly zero beyond its reach. \textbf{That is the dissociation.}
The model sees the context; the query that reads its state does not.

\textbf{An incidental finding, reported with its uncertainty.} We had predicted movement at $d=3$,
since the kernel is $4$. It is exactly zero. The reason is that \textbf{the oldest tap of the
convolution is exactly zero in all $24$ layers}, which we confirmed both from the weights and by
intervention: changing one token moves the convolution output at exactly three consecutive positions,
not four. So the effective window in this checkpoint is $3$, not $4$, and the reach is $2$. We do not
know the cause. A weight that is exactly zero across $24 \times 1536$ entries does not come from a
gradient, so a conversion artefact is plausible, but we did not verify that and we do not claim it.
The fact is reproducible in three lines. Its practical consequence, if it generalises, is that the
nominal kernel is an upper bound on the reach and should be measured rather than assumed.

\section{The diagnosis, as a recipe}

The measurement in Result 1 requires no training and no gradient. For any model that consults a memory
with a query built from a local window:

\begin{enumerate}
\item Take a query whose parts are separated, so that some component sits at distance $d$ from the read
position.
\item Replace that component with a different one of the same type.
\item Compare the retrieval distributions.
\end{enumerate}

If they are identical, that component is invisible to the search, and the model's confidence about it
is unearned. With a kernel of $4$, the default in deployed linear-attention layers, the reach is three
tokens, so any part of a question further back than that does not enter that layer's query.

\section{Limitations}

The substrate is a small model on a closed synthetic language, and the distance between the relation
and the read position is fixed by the generator, so the effect here is deterministic where in natural
text it would be a distribution. The measurement is of \emph{access}, not use: a token inside the window
is not necessarily used, and indeed the article tap is inside and contributes little. Finally, the read
in our model is injected before the mixer, which is what makes the zero exact; in a deep model that
consults memory from later layers, recurrence has already moved information forward and the effective
reach grows with depth. The claim is made for memory consulted from early layers, which is where it
lives.

\begin{thebibliography}{99}
\bibitem[Dao and Gu(2024)]{mamba2} Dao, T., and Gu, A. (2024). Transformers are SSMs: Generalized Models and Efficient Algorithms Through Structured State Space Duality. ICML 2024.
\bibitem[Yang et al.(2025)]{gateddeltanet} Yang, S., Kautz, J., and Hatamizadeh, A. (2025). Gated Delta Networks: Improving Mamba2 with Delta Rule. ICLR 2025. arXiv:2412.06464.
\bibitem[Hugging Face(2026)]{qwen3next} Hugging Face Transformers documentation, \texttt{Qwen3NextConfig}, parameter \texttt{linear\_conv\_kernel\_dim}, default $4$. Accessed 2 September 2026.
\bibitem[Li et al.(2024)]{cat} Li, M., Zhang, X., Huang, Y., and Oymak, S. (2024). On the Power of Convolution Augmented Transformer. arXiv:2407.05591.
\bibitem[Arora et al.(2024)]{zoology} Arora, S., Eyuboglu, S., Timalsina, A., Johnson, I., Poli, M., Zou, J., Rudra, A., and R\'e, C. (2024). Zoology: Measuring and Improving Recall in Efficient Language Models. ICLR 2024. arXiv:2312.04927.
\bibitem[Golovneva et al.(2025)]{mta} Golovneva, O., Wang, T., Weston, J., and Sukhbaatar, S. (2025). Multi-Token Attention. arXiv:2504.00927.
\end{thebibliography}

\end{document}
